\phantomsection
\addcontentsline{toc}{chapter}{TÓM TẮT}

\begin{center}
{\large\bfseries{TÓM TẮT}}
\end{center}

\medskip

\indent Video ngắn đang trở thành một dạng nội dung phổ biến, nhưng việc phân tích tự động loại video này vẫn gặp nhiều khó khăn vì nhịp dựng nhanh, cảnh quay thay đổi liên tục và hiệu ứng chuyển cảnh ngày càng đa dạng. Trong bối cảnh đó, phát hiện ranh giới cảnh quay là bước quan trọng để hệ thống có thể hiểu cấu trúc video trước khi thực hiện các tác vụ như tóm tắt, lập chỉ mục hoặc truy hồi nội dung.

\indent Nghiên cứu này tập trung vào bài toán phát hiện ranh giới cảnh quay cho video ngắn. Trước hết, các hướng tiếp cận tiêu biểu được tổng hợp để làm rõ cách bài toán đã được xử lý từ các phương pháp dựa trên đặc trưng thủ công đến các mô hình học sâu. Từ nền tảng đó, mô hình AutoShot được chọn làm cơ sở để phân tích và cải tiến, với trọng tâm là điều chỉnh tầng phân loại và cách xử lý xác suất đầu ra nhằm giúp mô hình nhận biết ranh giới cảnh quay ổn định hơn.

\indent Đóng góp chính của đề tài là đưa ra một hướng cải tiến phù hợp với đặc điểm của video ngắn mà không cần thay đổi toàn bộ kiến trúc mô hình ban đầu. Cách kết hợp một mạng xương sống đã được tối ưu với các điều chỉnh ở tầng phân loại cho thấy tiềm năng hỗ trợ tốt cho bài toán phát hiện ranh giới cảnh quay, đồng thời cũng đặt ra nhu cầu tiếp tục nâng cao khả năng tổng quát hóa khi áp dụng sang các miền dữ liệu khác. Những phân tích này có thể làm cơ sở cho các nghiên cứu tiếp theo về phân đoạn video và xây dựng dữ liệu video ngắn tiếng Việt.
